% main.tex -- schelet de articol A1 (C1: benchmark middleware ROS 2 sub degradare)
% LOCAL: \documentclass{article} + \bibliographystyle{plain} (IEEEtran lipseste local).
% OVERLEAF / TeX Live complet: comuta pe \documentclass[conference]{IEEEtran}.
% TOATE cifrele de mai jos sunt PROVIZORII (SIL = loopback, N=10) -- de inlocuit cu HIL.
% ASCII 100%. NU contine citari inventate: referintele sunt marcate TODO in references.bib.
\documentclass[11pt]{article}
\usepackage{graphicx}
\usepackage{booktabs}
\usepackage{amsmath}
\usepackage[margin=1in]{geometry}
\graphicspath{{figs/}}

\title{Selectia middleware-ului ROS 2 sub degradare controlata de retea:\\
       un benchmark CycloneDDS vs Zenoh si limitele unui selector invatat\\
       \large (DRAFT -- rezultate PROVIZORII pe SIL/loopback, N=10)}
\author{Alexandru Gheorghita \\ (afilliere TODO)}
\date{DRAFT \today}

\begin{document}
\maketitle

\begin{abstract}
% TODO proza finala. Schelet onest:
Teleoperarea robotilor peste retele degradate depinde de stratul de middleware (rmw)
din ROS 2. Comparam doua implementari, CycloneDDS si Zenoh, sub degradare controlata cu
\texttt{tc netem} (pierdere, latenta, jitter) pe trei dimensiuni de sarcina utila, si
intrebam daca un selector invatat poate alege per-conditie middleware-ul optim.
Pe N=10 repetitii in regim SIL (loopback, o singura masina), CycloneDDS ofera o coada de
latenta (RTT p95) mai mica si mai predictibila; un selector invatat, validat
leave-one-condition-out, NU bate regula triviala always-CycloneDDS pentru un obiectiv de
latenta. Sub un obiectiv constient de pierdere, politica optima devine DEPENDENTA DE
DEADLINE. Rezultatele sunt provizorii (loopback); comparatia autoritara cere HIL (doua
masini + legatura reala), pentru care furnizam un runbook.
\end{abstract}

\section{Introducere}
Teleoperarea robotilor in misiuni de cautare-salvare (SAR) si in reabilitare se desfasoara
peste legaturi wireless care se degradeaza imprevizibil: pierderea de pachete, latenta si
jitter-ul cresc cu distanta, cu obstacolele si cu interferenta. In ROS 2, comportamentul sub
aceasta degradare depinde de stratul de middleware (rmw) si de stiva de transport (DDS sau
Zenoh), nu doar de codul aplicatiei; alegerea middleware-ului devine astfel o decizie de
proiectare cu efect direct asupra calitatii teleoperarii.

Intrebarea practica e: exista un middleware uniform mai bun, sau alegerea optima depinde de
conditia de retea si de sarcina utila -- si, daca da, poate fi ea \emph{invatata}? Raspundem
empiric, comparand CycloneDDS si Zenoh sub degradare controlata cu \texttt{tc netem} si
incercand sa construim un selector per-conditie, evaluat onest.

Contributii: (i) un benchmark reproductibil CycloneDDS vs Zenoh sub \texttt{tc netem};
(ii) un cadru de \emph{selectie} a middleware-ului cu validare leave-one-condition-out
(LOCO) si metrica de \emph{regret}; (iii) un rezultat negativ onest si conditiile sub care
un selector ar deveni justificat; (iv) un drum SIL$\rightarrow$HIL identic ca schema de date.

\section{Lucrari conexe}
% TODO: DDS/QoS, Zenoh, emulare de retea (netem), selectie de middleware. Vezi references.bib.

\section{Metodologie}
\subsection{Nucleu pur + selftest}
Algoritmii (statistici RTT, features de conditie, regret, validare LOCO, cost constient de
pierdere) traiesc intr-un modul fara ROS, testat in izolare
(\texttt{selector\_core.py}, \texttt{bench\_core.py}; 30 respectiv 12 verificari).
\subsection{Degradarea retelei}
Conditiile sunt parametri \texttt{netem} reali (pierdere \%, latenta ms, jitter ms),
aplicati pe interfata. In SIL, interfata e \texttt{lo} (loopback); in HIL, interfata fizica.
\subsection{Selectorul si obiectivele}
Pentru fiecare celula (conditie $\times$ payload) alegem rmw-ul care minimizeaza un obiectiv:
\begin{itemize}
  \item \textbf{CONTROL}: $\min$ RTT p95.
  \item \textbf{CONSTIENT DE PIERDERE}: $\min$ cost asteptat
        $(1-\text{loss})\cdot \text{RTT}_{p95} + \text{loss}\cdot D$, unde $D$ e un deadline
        de control (cost atribuit unui esantion pierdut).
  \item \textbf{TELEMETRIE}: $\min$ timp de misiune (necesita stratul de misiune; groundwork
        construit, date reale in curs).
\end{itemize}
Validarea e leave-one-condition-out (NU split aleator: repetitiile aceleiasi celule ar
scurge informatie intre train si test). Raportam REGRET fata de oracol si fata de regulile
globale always-CycloneDDS / always-Zenoh.

\section{Configuratie experimentala}
% TODO: hardware, versiuni (ROS 2 Jazzy), payload-uri (64/4096/65536 B), N, conditii.
SIL: o singura masina, transport prin \texttt{lo}, \texttt{tc netem} local, N=10 repetitii,
8 conditii, 3 payload-uri. Memoria partajata DDS dezactivata; descoperire P2P.
\textbf{PROVIZORIU}: toate cifrele de mai jos sunt SIL; HIL ramane standardul autoritar
(vezi \texttt{HIL\_RUNBOOK.md}).

\section{Rezultate (PROVIZORII -- SIL, N=10)}
\subsection{Transport: latenta de coada}
% Figura: fig_transport (RTT p95 per conditie, grupat pe RMW).
\begin{figure}[t]\centering
  % \includegraphics[width=\linewidth]{fig_transport} % TODO: copiaza figura in figs/
  \fbox{\texttt{figs/fig\_transport.pdf (TODO: copiaza din rularea campaniei)}}
  \caption{RTT p95 per conditie, CycloneDDS vs Zenoh (SIL, N=10). PROVIZORIU.}
  \label{fig:transport}
\end{figure}
Pe cele 24 de celule, CycloneDDS castiga 17/24; always-CycloneDDS are regret mediu
\textbf{142 ms} vs always-Zenoh \textbf{1143 ms}. CycloneDDS e si mai predictibil
(coeficient de variatie al p95 mai mic). [Cifre PROVIZORII, SIL.]

\subsection{Selectorul de control nu bate regula triviala}
% Figura: selector_regret.
\begin{figure}[t]\centering
  \fbox{\texttt{figs/selector\_regret.pdf (TODO)}}
  \caption{Regret mediu (LOCO) pe strategie de selectie, obiectiv control. PROVIZORIU.}
  \label{fig:selector}
\end{figure}
Sub validare LOCO, selectorul invatat (1-NN: 66.7\% acuratete, regret 558 ms; arbore: 653 ms)
NU bate always-CycloneDDS (142 ms). \textbf{Rezultat negativ onest}: datele disponibile nu
justifica un selector per-conditie pentru obiectivul de latenta.

\subsection{Obiectiv constient de pierdere: politica dependenta de deadline}
Sub costul asteptat cu deadline $D$, gapul se ingusteaza, iar la $D=5000$ ms selectorul 1-NN
(\textbf{104 ms}) BATE always-CycloneDDS (\textbf{188 ms}). Concluzie: un selector isi merita
locul doar in regimul unde un esantion pierdut e foarte scump (control de siguranta).
[Sensibilitate la $D$; PROVIZORIU.]

\subsection{Caracterizare ML (PDIA)}
Un model descriptiv (Ridge pe severitate + log payload) prezice log RTT cu $R^2_{test}=0.90$;
un clasificator de "link utilizabil" ($<500$ ms) atinge acuratete $0.94$. [PROVIZORIU.]

\section{Discutie}
Rezultatele trebuie citite prin limitele lor. (1) Sunt SIL: transportul trece prin
\texttt{lo} (loopback) pe o singura masina, ceea ce subevalueaza diferentele dintre stivele
reale -- de aici necesitatea HIL. (2) Sub \texttt{netem}, conditiile noastre de "pierdere"
modeleaza legaturi fara retransmisie (UDP brut); pe Wi-Fi real, stratul MAC retransmite
singur, deci pierderea apare mai degraba ca latenta si jitter. (3) Burstiness-ul (corelatia
pierderilor) nu este inca un feature al selectorului: conditiile \texttt{*\_burst} sunt
tratate cu aceiasi parametri ca varianta de baza. (4) Obiectivul de telemetrie cere stratul
de misiune (timpi de misiune per rmw); infrastructura e construita, datele reale sunt in curs.
Aceste limite indica direct campania HIL si campania de misiune ca pasi urmatori.

% TODO: Lucrari conexe (cu citari REALE in references.bib) -- nu se inventeaza aici.

\section{Concluzii si lucrari viitoare}
Pe SIL, CycloneDDS domina pentru control si un selector nu se justifica; sub un obiectiv
constient de pierdere, justificarea devine dependenta de deadline. \textbf{Urmatorul pas:}
campanie HIL pe doua masini (PC + RaspberryPi 4B) cu schema de date identica, plus stratul de
misiune pentru obiectivul telemetrie. Vezi \texttt{HIL\_RUNBOOK.md}.

\bibliographystyle{plain}
\bibliography{references}

\end{document}
